\subsection{Local construction of reflexive modules}

We keep the notation and hypotheses of no. 2. We shall say that a property holds ``for almost all \(p \in P\) '' if the set of \(p \in P\) for which it is not true is finite.

THEOREM 3. Let V be a vector space of finite rank over K and M a lattice of V with respect to A.

\begin{enumerate}
\def\labelenumi{(\roman{enumi})}
\item
  Let N be a lattice of V with respect to A; then, for every prime ideal p of A, \(N_{p}\) is a lattice of V with respect to \(A_{p}\) and, for almost all \(p \in P\) , \(N_{p} = M_{p}\) .
\item
  Conversely, suppose given for all \(\mathfrak{p} \in \mathbf{P}\) a lattice \(\mathrm{N}(\mathfrak{p})\) of \(\mathbf{V}\) with respect to \(\mathbf{A}_{\mathfrak{p}}\) such that \(\mathrm{N}(\mathfrak{p}) = \mathrm{M}_{\mathfrak{p}}\) for almost all \(\mathfrak{p} \in \mathbf{P}\) . Then \(\mathrm{N} = \bigcap_{\mathfrak{p} \in \mathbf{P}} \mathrm{N}(\mathfrak{p})\) is a reflexive lattice of \(\mathbf{A}\) with respect to \(\mathbf{A}\) and it is the only reflexive lattice \(\mathrm{N}'\) of \(\mathbf{V}\) with respect to \(\mathbf{A}\) such that \(\mathrm{N}_{\mathfrak{p}}' = \mathrm{N}(\mathfrak{p})\) for all \(\mathfrak{p} \in \mathbf{P}\) .
\item
  The first assertion follows from no. 1, Proposition 4. Moreover, there exist \(x, y\) in \(\mathbf{K}^*\) such that \(x\mathrm{N} \subset \mathrm{M} \subset y\mathrm{N}\) (no. 1, Proposition 2); we know that, for almost all \(\mathfrak{p} \in \mathbf{P}\) , \(v_{\mathfrak{p}}(x) = v_{\mathfrak{p}}(y) = 0\) (§ 1, no. 6, Theorem 4), which shows that \(x\) and \(y\) are invertible in \(\mathbf{A}_{\mathfrak{p}}\) and hence \(\mathbf{M}_{\mathfrak{p}} = \mathbf{N}_{\mathfrak{p}}\) .
\item
  We may replace M by \(x^{-1}\mathbf{M}\) where \(x \neq 0\) in A and assume that \(\mathrm{N}(\mathfrak{p}) \subset \mathbf{M}_{\mathfrak{p}}\) for all \(\mathfrak{p} \in \mathrm{P}\) . Let \(\mathfrak{p}_1, \ldots, \mathfrak{p}_h\) be the elements of P such that \(\mathrm{N}(\mathfrak{p}) = \mathrm{M}_{\mathfrak{p}}\) for \(\mathfrak{p}\) distinct from the \(\mathfrak{p}_i\) ( \(1 \leqslant i \leqslant h\) ); we write:
\end{enumerate}

\[
Q = M \cap N (p _ {1}) \cap \dots \cap N (p _ {h}).
\]

As each of the \(\mathrm{N}(p_{i})\) contain a free lattice with respect to \(A_{p_{i}}\) , it contains a fortiori a lattice of V with respect to A, hence Q contains a lattice of V with respect to A (no. 1, Proposition 3) and, as Q is contained in M, Q is a lattice with respect to A. To prove that \(Q_{p} = N(p)\) for all \(p \in P\) , we shall use the following lemma:

LEMMA 1. Let \(\mathfrak{p}\) and \(\mathfrak{p}'\) be two prime ideals of \(\mathbf{A}\) such that (0) is the only prime ideal of \(\mathbf{A}\) contained in \(\mathfrak{p} \cap \mathfrak{p}'\) . Then, for every sub- \(\mathbf{A}\) -module \(\mathbf{E}\) of \(\mathbf{V}\) , \((\mathbf{E}_{\mathfrak{p}})_{\mathfrak{p}'} = \mathbf{K} \cdot \mathbf{E}\) .

Let S be the multiplicative subset \((\mathbf{A}-\mathfrak{p})(\mathbf{A}-\mathfrak{p}')\) of A; by Chapter II, §2, no. 3, Proposition 7, \((\mathrm{E}_{\mathfrak{p}})_{\mathfrak{p}'}=\mathrm{S}^{-1}\mathrm{E}\) . Further, \(A\subset S^{-1}A\subset K\) ; the prime ideals of \(S^{-1}A\) correspond to the prime ideals q of A such that \(q\cap S=\varnothing\) (Chapter II, §2, no. 5, Proposition 11) and by hypothesis (0) is the only prime ideal of A not meeting S; hence \(S^{-1}A=K\) and \(S^{-1}E=K.E\) .

We now return to the proof of (ii). If \(p \in P\) is distinct from the \(p_i (1 \leqslant i \leqslant h)\) , Lemma 1 applied to \(\mathrm{N}(\mathfrak{p}_i)\) gives \((\mathrm{N}(\mathfrak{p}_i))_{\mathfrak{p}} = ((\mathrm{N}(\mathfrak{p}_i))_{\mathfrak{p}_i})_{\mathfrak{p}} = \mathrm{K}. \mathrm{N}(\mathfrak{p}_i) = \mathrm{V}\) , since the \(p_i\) and p are of height 1. Then

\[
\mathbf {Q} _ {\mathfrak {p}} = \mathbf {M} _ {\mathfrak {p}} \cap (\mathbf {N} (\mathfrak {p} _ {1})) _ {\mathfrak {p}} \cap \dots \cap (\mathbf {N} (\mathfrak {p} _ {h})) _ {\mathfrak {p}} = \mathbf {M} _ {\mathfrak {p}} = \mathbf {N} (\mathfrak {p})
\]

(Chapter II, § 2, no. 4). On the other hand, if \(\mathfrak{p}\) is equal to \(\mathfrak{p}_i\) ( \(1 \leqslant i \leqslant h\) ), then \((\mathrm{N}(\mathfrak{p}_i))_{\mathfrak{p}_j} = \mathrm{V}\) for \(i \neq j\) by the argument as above and \((\mathrm{N}(\mathfrak{p}_i))_{\mathfrak{p}_i} = \mathrm{N}(\mathfrak{p}_i)\) , whence

\[
\mathbf {Q} _ {\mathfrak {p} _ {i}} = \mathbf {M} _ {\mathfrak {p} _ {i}} \cap \mathbf {N} (\mathfrak {p} _ {i}) = \mathbf {N} (\mathfrak {p} _ {i}).
\]

We have therefore proved that \(Q_{\mathfrak{p}} = N(\mathfrak{p})\) for all \(\mathfrak{p} \in P\) . Then

\[
\mathbf {N} = \mathbf {Q} ^ {* *} = \bigcap_ {\mathfrak {p} \in \mathbf {P}} \mathbf {Q} _ {\mathfrak {p}}
\]

is reflexive and satisfies the relations \(N_{p} = Q_{p} = N(p)\) for all \(p \in P\) ; the uniqueness property follows immediately from Theorem 2 of no. 2.

Remark. Let L be a free lattice of V with respect to A. Since \(A_{p}\) is a principal ideal domain for \(p \in P\) , \(N(p)\) is a free \(A_{p}\) -module of the same rank as L and there exists \(u(p) \in \mathbf{GL}(V)\) such that \(u(p)(L_{p}) = N_{p}\) ; this condition moreover determines \(u(p)\) to within right multiplication by an element of \(\mathbf{GL}(L_{p})\) . The condition \(N(p) = L_{p}\) for almost all \(p \in P\) means that necessarily \(u(p) \in \mathbf{GL}(L_{p})\)

for almost all \(\mathfrak{p} \in \mathbb{P}\) . The families \((u(\mathfrak{p}))_{\mathfrak{p} \in \mathbb{P}}\) satisfying the latter property form a multiplicative group \(\mathbf{GL}_a(V)\) containing as subgroup the product \(\prod_{\mathfrak{p} \in \mathbb{P}} \mathbf{GL}(L_{\mathfrak{p}})\) . Theorem 3 then shows that the set of reflexive lattices of \(V\) is in canonical one-to-one correspondence with the homogeneous space \(\mathbf{GL}_a(V)/\prod_{\mathfrak{p} \in \mathbb{P}} \mathbf{GL}(L_{\mathfrak{p}})\) . If a basis \((e_i)_{1 \leq i \leq n}\) of \(L\) over \(A\) is chosen, \(\mathbf{GL}(V)\) (resp. \(\mathbf{GL}(L_{\mathfrak{p}})\) ) is identified with the group of invertible matrices \(\mathbf{GL}(n, K)\) (resp. \(\mathbf{GL}(n, A_{\mathfrak{p}})\) ) and the group \(\mathbf{GL}_a(V)\) with the group of systems of matrices of order \(n\) , \((U(\mathfrak{p}))_{\mathfrak{p} \in P}\) , such that \(U(\mathfrak{p}) \in \mathbf{GL}(n, K)\) for all \(\mathfrak{p} \in P\) and \(U(\mathfrak{p}) \in \mathbf{GL}(n, A_{\mathfrak{p}})\) for almost all \(\mathfrak{p} \in P\) . If \(A\) is a Dedekind domain, the group \(\mathbf{GL}_a(V)\) is also identified with the group \(\mathbf{GL}(n, A)\) , where \(A\) is the ring of restricted adèles (§ 2, no. 4).

